\section{Canonical quotients and the complete prime-to-$p$ index}\label{sec:canonical-iteration}
This section proves the complete index table by iterating central rigidity while retaining the actual torsion comparison at every quotient.

Throughout, $p\ge5$, $K$ is complete discretely valued of characteristic zero with perfect residue field, $\mu_p\subset K$, and $p\nmid e=e_K$. For potentially supersingular groups use the common tame good field $L/K$ of degree $p^2-1$ and the normalized Hasse parameter of Section~\ref{sec:central-ramification}. The inputs are the pointwise graph bound, the one-slope and prime-Hasse recovery statements, the actual quotient formulas, and central $p^2$ rigidity, each proved in the preceding sections.

\begin{lemma}[Iteration with transported kernels]\label{ci-lem:iteration}
Let $f:G[p^n]\simeq H[p^n]$ be an isomorphism of potentially supersingular height-two groups. Suppose $G$ has a two-slope parameter $z$ divisible by $p$. If an integer $c$ satisfies
\[
 0\le c\le n-2,\qquad e<p^c z<pe,
 \qquad p^jz<e\quad(0\le j<c),
\]
then $\theta_G=\theta_H$.
\end{lemma}

\begin{proof}
Put $G_0=G$, $H_0=H$ and $f_0=f$. For $j<c$, let $C_j$ be the canonical $p$-line of $G_j$. It is $K$-stable: the semilinear good-model action preserves formal valuations, and the high line consists of zero and all points of the larger first-level valuation. Define
\[
 D_j=f_j(C_j),\qquad
 G_{j+1}=G_j/C_j,\qquad H_{j+1}=H_j/D_j.
\]
The subgroup $D_j$ is $K$-stable because $f_j$ is Galois-equivariant. Lemma~\ref{lem:actual-quotients} gives the actual quotients over $K$ and preserves their special characters. No condition that $D_j$ be canonical in $H_j$ is required.

To check the retained truncation precisely, put $q=n-j$. As finite Galois modules,
\begin{equation}\label{ci-eq:transport}
 (G_j/C_j)[p^{q-1}]
 =\{x\in G_j[p^q]:p^{q-1}x\in C_j\}/C_j.
\end{equation}
Indeed a quotient point killed by $p^{q-1}$ has a lift $x$ with $p^{q-1}x\in C_j$, hence $p^qx=0$ since $pC_j=0$; the converse follows by multiplication. The map $f_j$ carries this exact subquotient to its partner for $H_j,D_j$. It therefore defines an isomorphism $f_{j+1}$ at level $p^{n-j-1}$. Each step loses exactly one level and requires no integral extension of the original torsion isomorphism.

For every step taken, $p^jz<e$. The canonical quotient formula of Proposition~\ref{rc-prop:quotients} gives the parameter $z_j=p^jz$ of $G_j$ by induction. The partner remains potentially supersingular because it is $K$-isogenous to $H$. At the final step, $n-c\ge2$ and
\[
 e<z_c<pe,\qquad p\mid z_c,
 \qquad p\nmid (p+1)e-z_c.
\]
The last assertion uses $p\nmid e$. Restrict $f_c$ to $p^2$ and apply Theorem~\ref{thm:central-p2-rigidity}. The terminal traces agree, and all intervening $K$-isogenies preserve those traces. Hence the original traces agree.
\end{proof}

\begin{proof}[Proof of Theorem~\ref{thm:ss-iterated-range}]
Suppose the special characters differ. Since $n\ge3$, Lemma~\ref{lem:one-slope-recovery} excludes every case involving a one-slope model. Both models are therefore two-slope. If both normalized Hasse values were prime to $p$, Theorem~\ref{rc-thm:primehasse} would give equal traces already from $p$-torsion. Choose the model $G$ so that its parameter $z$ is divisible by $p$.

The smallest nonzero $p^n$-torsion value of $G$ is $z/p^{2n-1}$. To recall the range of the division calculation, its first-level minimum is $B=z/p$. A nonzero torsion target of value $q\le A=(p+1)e-z$ has all $p$-division values $q/p^2$ when $q\le p^2B$; otherwise the division values are $(q-p(p-1)B)/p$ and $B$. Multiplication by $p$ strictly increases positive nonzero valuations, so every target stays within $q\le A$. The smallest targets always divide by $p^2$; larger targets below $p^2B$ yield larger divisions, and targets above that threshold give divisions at least $B$. Division points of zero are first-level points and do not lower the minimum. Induction proves the asserted minimum, as also established before~\eqref{cr-eq:zbound}.

The pointwise graph bound, Proposition~\ref{fc-thm:valuation}, now gives
\[
 z\ge Z_n:=\frac{(p-1)p^{2n-1}}{p^n-1}.
\]
The assumed bound on $e$ is $e<p^{n-2}Z_n$, so $p^{n-2}z>e$. Let $c$ be the least nonnegative integer with $p^cz>e$. It exists and $c\le n-2$. Equality $p^jz=e$ is impossible because $p\mid z$ and $p\nmid e$. If $c=0$, the two-slope range gives $z<pe$. If $c>0$, minimality gives $p^{c-1}z<e$, hence $p^cz<pe$. Lemma~\ref{ci-lem:iteration} applies and contradicts the unequal characters.

For the integral form of the bound, use
\[
 \frac{p^{3n-3}}{p^n-1}
 =p^{2n-3}+p^{n-3}+\frac{p^{n-3}}{p^n-1}.
\]
The final term is strictly between zero and one, and $e/(p-1)$ is an integer.
\end{proof}

\begin{proof}[Proof of Corollary~\ref{cor:complete-prime-recovery}]
The supersingular bound is larger than the asserted full-class threshold, since
\[
 \frac{(p-1)p^{3n-3}}{(p^n-1)p^{2n-3}}
 =\frac{(p-1)p^n}{p^n-1}>1.
\]
We give the structural reduction to make its lack of an additional upper bound on $e$ explicit. Determinants have the form $\chi_p^d\epsilon$, with $d\in\{0,1,2\}$ and $\epsilon$ of order dividing $p-1$. Because $p\nmid e$, the wild cyclotomic order modulo $p^n$ is $p^{n-1}$. Equality of determinant truncations forces $d=d'$, since $|d-d'|\le2<p$, and then equality of the finite factors.

In dimension zero, Brauer lifting of the semisimplified residual character recovers the special character; Cartier duality gives dimension two. In dimension one, the connected--\'etale filtration of a potentially ordinary piece descends by uniqueness. The possibilities are therefore a whole supersingular group or two height-one graded pieces. For two filtered groups, $\chi_p\equiv1\pmod p$ identifies their special traces with the Brauer lifts of their residual semisimplifications. A mixed comparison would require $e>p^{2n-3}$ by Proposition~\ref{prop:mixed-obstruction}, contrary to the hypothesis. The remaining supersingular comparison is covered by Theorem~\ref{thm:ss-iterated-range}. This proves the uniform full-class bound.
\end{proof}

\begin{proof}[Proof of Theorem~\ref{thm:complete-prime-index}]
At $e=p-1$, Theorem~\ref{thm:minimal-index} gives the exact index one and an actual pair witnessing failure at the first level. Every other positive multiple of $p-1$ is at least $2(p-1)>p$. Since $p\nmid e$, there is a unique $r\ge2$ such that
\[
 p^{2r-3}<e<p^{2r-1}.
\]
Corollary~\ref{cor:complete-prime-recovery} with $n=r+1$ gives $\nuMin(p,e,2)\le r$. Corollary~\ref{cor:sharp-mixed-prime} at level $p^r$ gives an actual unequal-character comparison over some field of index $e$, because $e>p^{2r-3}$. Thus $\nuMin(p,e,2)\ge r$. Both inequalities are attained at every permitted index, proving the exact table. Solving its strict inequalities yields~\eqref{eq:complete-prime-index}; no equality endpoint occurs because $p\nmid e$.
\end{proof}

For example, at $p=5$ the index is three for every allowed $125<e<3125$, and four for every allowed $3125<e<78125$. Here ``allowed'' means $4\mid e$, $5\nmid e$, with $\mu_5$ in the base. The lower constructions vary the field and are actual Barsotti--Tate comparisons; the upper assertions are uniform over every field in the specified class.
